\par\Needspace{8\baselineskip}\bigskip
\pdfbookmark[3]{Exposition, 16}{pallas-II-article16}
\begin{center}16.\end{center}
\noindent\hspace*{\parindent}Si au lieu de \(T,V\) on préfère de faire usage de \(R,S\), la forme la plus commode du calcul paraît être de rapporter le lieu tant de la planète perturbante que de la planète troublée à trois plans perpendiculaires entre eux, dont l’un soit le plan même de l’orbite osculatrice de la dernière planète, l’un des autres passant par la ligne d’apsides de cette planète. Soient \(X',Y',Z'\) les coordonnées de la planète perturbante, \(X,Y,0\) celles de la planète troublée, savoir
\SourceLeaf{466}
\[
\begin{aligned}
X'&=r'\cos PL',&\qquad X&=r\cos(v-\tilde\omega)\\
Y'&=r'\cos\Pi L',&Y&=r\sin(v-\tilde\omega)\\
Z'&=r'\cos ML',&&
\end{aligned}
\]
on aura
\[
\begin{gathered}
\rho=\BareRoot{((X'-X)^2+(Y'-Y)^2+Z'Z')}\\[4pt]
\begin{aligned}
\dfrac{R}{\mu}&=\dfrac{X'-X}{\rho^3}-\dfrac{X'}{r'^2}\\
\dfrac{S}{\mu}&=\dfrac{Y'-Y}{\rho^3}-\dfrac{Y'}{r'^3}\\
\dfrac{W}{\mu}&=\dfrac{Z'}{\rho^3}-\dfrac{Z'}{r'^3}\mathpunct{\cdot}
\end{aligned}
\end{gathered}
\]
\par La manière la plus commode de calculer ces coordonnées étant suffisamment connue, il serait superflu de nous y arrêter.
\par Puisque dans toutes les recherches précédentes, les différens arcs sont toujours censés être exprimés en parties du rayon, il est clair que dans toutes les formules homogènes, c. à d. où les arcs ont même dimension des deux côtés, on peut aussi supposer les arcs exprimés en secondes, mais que dans celles, qui ne le sont pas, il faut ajouter le facteur 206265 du côté où les arcs ont une dimension de moins. Ainsi les équations 61, 62, 66, 68, 70, 74, 75, 76 resteront les mêmes, mais dans les équations 63, 64, 65, 72, 73 il faudra diviser la seconde partie par 206265, si tous les arcs sont censés exprimés en secondes. Au contraire, si l’on voulait en excepter l’arc \(n\), et qu’on continuait de l’exprimer en parties du rayon, on devrait multiplier la seconde partie des équations 61, 62, 66 etc. par 206265 tandisque les autres équations 63, 64, 65 etc. resteraient sans changement.
\par\medskip\begin{center}\rule{4em}{.4pt}\end{center}
